% Generated draft: requires mathematical and source review.
\documentclass[11pt]{article}
\usepackage[margin=25mm]{geometry}
\usepackage{fontspec}
\setmainfont{Latin Modern Roman}
\usepackage{amsmath,amssymb,mathtools,unicode-math}
\setmathfont{Latin Modern Math}
\usepackage{enumitem,needspace,xurl,hyperref}
\setlength{\parindent}{0pt}
\setlength{\parskip}{5pt}
\setlength{\emergencystretch}{3em}
\newcommand{\N}{\mathbb{N}}
\newcommand{\Z}{\mathbb{Z}}
\newcommand{\Q}{\mathbb{Q}}
\newcommand{\R}{\mathbb{R}}
\newcommand{\C}{\mathbb{C}}
\newcommand{\F}{\mathbb{F}}
\newcommand{\A}{\mathbb{A}}
\newcommand{\PP}{\mathbb{P}}
\newcommand{\E}{\mathbb{E}}
\newcommand{\eps}{\varepsilon}
\newcommand{\dd}{\,\mathrm{d}}
\begin{document}
\section*{UnsolvedMath ID: 30003274}
\textbf{Source problem code:} catalogue entry
\subsection{Definitions and mathematical statement}
Let \(q=p^r\), where \(p\) is prime and \(r\geq 1\), and let \(\F_q\) be the field with \(q\) elements. Let
\[
f(x)=ax^2+bx+c\in\F_q[x],\qquad a\neq 0,
\]
and let \(m\in\F_q\). Define \(x_0=m\) and \(x_{n+1}=f(x_n)\) for every integer \(n\geq 0\); thus \(x_n=f^{\circ n}(m)\), with \(f^{\circ 0}\) the identity map. Since \(\F_q\) is finite, this sequence is eventually periodic. Define its preperiod by
\[
\mu(f,m)=\min\bigl\{u\in\Z_{\geq 0}:\text{ there exists }v\in\Z_{\geq 1}\text{ such that }x_{u+v}=x_u\bigr\},
\]
and its eventual period by
\[
\lambda(f,m)=\min\bigl\{v\in\Z_{\geq 1}:x_{\mu(f,m)+v}=x_{\mu(f,m)}\bigr\}.
\]
The eventual cycle consists of the \(\lambda(f,m)\) distinct elements beginning at \(x_{\mu(f,m)}\), and \(x_{n+\lambda(f,m)}=x_n\) for every \(n\geq\mu(f,m)\). In particular, \(\mu(f,m)=0\) when the initial value already lies on a cycle.

\textbf{Problem.} For every prime power \(q\), every quadratic polynomial \(f\in\F_q[x]\), and every \(m\in\F_q\), determine the exact values of \(\mu(f,m)\) and \(\lambda(f,m)\) in terms of the given data.

\textbf{Unresolved qualification.} The supplied statement imposes no restriction on the characteristic or on the quadratic family. It does not specify whether the requested determination should be an explicit formula, a structural classification, or an algorithm with prescribed efficiency; in particular, it does not explain what is sought beyond direct iteration on the finite field.
\subsection{Short English statement}
Given a quadratic polynomial over a finite field and a starting value, determine how many iterations occur before the orbit first enters a cycle and how many distinct elements that cycle contains. The intended form or efficiency of a general determination is unspecified. The supplied registry labels the problem partially solved; this draft does not independently verify that assessment.
\subsection{Sources}
\begin{itemize}
\item Supplied local UnsolvedMath registry record, permanent ID \(30003274\), title \emph{Preperiods and Cycles of Quadratic Polynomial Iteration}, source problem code \texttt{catalogue entry}, in the Oberwolfach Reports Open Problems collection. The record gives proposed year 2016 and names no proposer.
\item Supplied workshop provenance: \emph{Analytic Number Theory (2017)}, Oberwolfach Reports, pp.\ 2975--3030, DOI \texttt{10.4171/owr/2016/53}; \url{https://doi.org/10.4171/owr/2016/53}. The supplied background describes these reports as documenting mathematical workshops and their open-problem discussions.
\item Supplied literature reference labeled ``Heath-Brown 2017'': \url{https://doi.org/10.1112/S0025579317000328}.
\item Supplied literature reference labeled ``Heath-Brown preprint'': \url{https://arxiv.org/abs/1701.02707}.
\item Literature assessment dated \texttt{2026-08-21}: supplied dated curation derived from the OWR record, not independently verified. It labels the problem partially solved and states that Heath-Brown proved strong cycle-length bounds and analyzed important quadratic families, while the exact cycle length for a general quadratic polynomial over a finite field remains undetermined. The supplied background explicitly requires independent verification of time-sensitive status and cited claims.
\item The registry labels the statement \texttt{corrected\_verified}, explaining that mathematical content was preserved while notation, typography, or omitted definitions were repaired for self-contained publication. This is supplied registry metadata, not independent verification of this draft.
\end{itemize}
\end{document}
